\renewcommand{\arraystretch}{1.25}
\begin{table}[H]
\centering
\begin{tabular}{|p{4cm}|p{7cm}|p{2cm}|}
\hline
\textbf{Aspek GAP Evaluasi Penelitian} & 
\textbf{Parameter Kebutuhan Fungsional yang Direpresentasikan} &
\textbf{Kode} \\ \hline

Representasi Struktur Data &
Sistem harus mengubah skema objek hasil parsing menjadi struktur representasi
yang terorganisir, membedakan setiap sumber daya, atribut, serta relasi referensial
(\texttt{\$ref}/\texttt{items.\$ref}). &
F-02 \\ \hline

Representasi Struktur Data &
Sistem harus menyimpan metadata atribut konfigurasi (deskripsi, tipe data,
dan status \textit{required/optional}) sebagai bagian integral dari representasi terstruktur,
agar hierarki konfigurasi tetap utuh dan tidak terputus. &
F-03 \\ \hline

Mekanisme Validasi &
Sistem harus menerapkan mekanisme pembatasan output agar hanya menghasilkan
YAML valid ketika pengguna meminta pembuatan konfigurasi, serta tidak menambahkan
field yang tidak tersedia dalam konteks dokumen API. &
F-07 \\ \hline

Ambiguitas Entitas dan Terminologi &
Sistem harus mampu mengidentifikasi entitas Kubernetes yang dimaksud oleh pengguna
dan membedakan tipe pertanyaan menjadi \textbf{Conceptual\_QA} atau
\textbf{Code\_Generation}, sehingga pemanggilan konteks dilakukan secara tepat
berdasarkan maksud pertanyaan. &
F-04 \\ \hline

Ambiguitas Entitas dan Terminologi &
Sistem harus menyediakan mekanisme retrieval untuk mengambil himpunan objek dan field
yang benar-benar relevan berdasarkan entitas target dan intensi pengguna,
menghindari pencampuran konteks antar objek akibat penggunaan keyword semata. &
F-05 \\ \hline

Penanganan Referensi Silang &
Sistem harus mengonversi referensi struktural (misalnya \texttt{\$ref}, pola selector,
atau hubungan template) menjadi representasi relasional eksplisit melalui node–edge
agar traversal logis dapat dilakukan secara deterministik. &
F-02 \\ \hline

Penanganan Referensi Silang &
Sistem harus menyusun hasil retrieval menjadi konteks teknis yang terstruktur
dan menyuntikkannya ke dalam \textit{prompt} sebagai landasan jawaban LLM,
agar keputusan model berbasis pada data API yang benar. &
F-06 \\ \hline

Pemeliharaan Pengetahuan &
Sistem harus membatasi ruang lingkup parsing hanya pada enam komponen target
(Deployment, Pod, Service, Ingress, ConfigMap, Secret),
agar pembaruan skema hanya memerlukan pergantian sumber data OpenAPI tanpa
penulisan ulang aturan manual. &
F-01 \\ \hline

Akurasi Jawaban &
Sistem harus mengekstrak skema terstruktur dan atribut kunci objek,
serta menggunakannya sebagai konteks pemodelan jawaban,
guna memastikan instruksi kode YAML yang dihasilkan memiliki struktur valid. &
F-06 \\ \hline

Akurasi Jawaban &
Sistem harus membatasi output YAML hanya ketika intent diklasifikasikan sebagai 
\textbf{Code\_Generation}, memastikan format jawaban berbeda untuk pertanyaan 
konseptual sehingga hasil terbaca sesuai ekspektasi pengguna. &
F-07 \\ \hline

\textbf{Konsistensi Penggunaan Sistem (Tambahan — Non-Structural GAP)} &
Sistem harus menyediakan antarmuka \textit{Command-Line Interface (CLI)} 
untuk menjalankan fungsi inti (\texttt{generate}, \texttt{explain},
\texttt{validate}, \texttt{--help}), sehingga interaksi dapat dilakukan secara
ringan, konsisten, dan tanpa ketergantungan komponen visual. &
F-08 \\ \hline

\end{tabular}
\caption{Pemetaan Aspek GAP Penelitian terhadap Kebutuhan Fungsional Sistem}
\label{tbl:gap_mapping_functional}
\end{table}
